\section{Frontier Models}
\sectionkicker{FRONTIER PRICING}
\begin{wideflow}
{\Large\bfseries 値札の並んだ最前線――Sonnet 5.5、GPT-6.1 Sol、Gemini 4 Argon}\par
\smallskip
{\color{SurveyMuted}9月28日から30日に三つの旗艦が並んだ。値段と提供の形、発表側の物差しを切り分けて読む。}\par
\end{wideflow}

9月28日、AnthropicはClaude Sonnet 5.5を出した\cite{sonnet55}。顔つきは堅実な改良で、番うのは値札と守りの姿勢である。腕前の物差しは発表側の手による。Terminal-Benchで70.6\%、GDPval-AAで1844と置かれるが、いずれも売り手の計り方として読む。本号では独立の再走はしていない。

9月29日、OpenAIはGPT-6.1 Solを出した\cite{gpt61sol}。値札は確かめが取れている。入力100万トークンあたり2ドル、出力10ドル、キャッシュ読みが0.10ドルで、変更履歴との突き合わせも済んでいる。速さをうたうUltrafastのsoonは今号の話に入れない。能力の物差しに出る差分は発表側の計り方であり、独立の再現待ちとして置く。日付は日の単位までしか確定しない。

9月30日、GoogleはGemini 4 Argonを出した\cite{geminiargon}。読み違えてはならないのは100万という枠で、これは出力トークンの上限であり、入力でも到達済みの推論の話でもない。導入時の値付けと売り手の到達点は帰属つきで添える。試用寄りの構えという位置づけも崩さない。三つ並べると、今週の最前線は頂の交代ではなく、値段と渡し方の競争だったことが見える。

\begin{claimboundary}[資料と評価の境界]
料金表は発表時点の公式ページであり、請求の実測ではない。評価の詳細は本号の検証範囲外であり、数値はベンダー側の報告として独立した再現を待つ。相手の数値は公開報告からの引用であり、横断的な順位づけには使わない。出力100万の枠を入力や到達済み推論と読み替えない。
\end{claimboundary}
